%&latex
\documentclass[a4paper]{article}

\frenchspacing

\usepackage[cp1250]{inputenc}
\usepackage[czech]{babel}

\usepackage{a4wide}
\usepackage{amsmath, amsthm, amssymb, amsfonts}
\usepackage[mathcal]{eucal}
\usepackage{graphicx}
\usepackage{url}
\usepackage{color}
\usepackage{wrapfig}
\usepackage{capt-of}
\usepackage{float}
\usepackage{ifthen}
\usepackage{subfig}
\usepackage[normalem]{ulem}


% sirka a vyska textu nastavena jako papir, vsechny okraje vynulovany a pridano 20pt na kazdou stranu
% horizontalni rozmery
\setlength{\textwidth}{\paperwidth}
\addtolength{\textwidth}{-40pt}
\addtolength{\hoffset}{-1in}
\addtolength{\hoffset}{20pt}
\setlength{\oddsidemargin}{0in}
\setlength{\marginparsep}{0in}
% vertikalni rozmery
\setlength{\textheight}{\paperheight}
\addtolength{\textheight}{-60pt}
\addtolength{\voffset}{-1in}
\addtolength{\voffset}{20pt}
\setlength{\topmargin}{0in}
\setlength{\headheight}{0in}
\setlength{\headsep}{0in}


%Obrazek na miste
%pouziti
%%\obrazeknahore{adresa}{popisek}{label}
\long\def\obrazeknahore#1#2#3 {

\begin{figure}[t]
    \centering
    \includegraphics[width=0.8\textwidth]{#1}
    
    \caption{#2}
    \label{#3}
    
\end{figure}

}


%==========================================
%PEKELNA MAKRA NA ZAROVNANI OBRAZKU DOPRAVA

\makeatletter


%tohle je makro, ktere mi dovoluje obtekani i u kratkych environmentu
%ABSOLUTNE nechapu, jak to funguje, ale funguje to
%viz http://tex.stackexchange.com/questions/26078/ 
\def\odrovnej{\@@par
\ifnum\@@parshape=\z@ \let\WF@pspars\@empty \fi % reset `parshape'
\global\advance\c@WF@wrappedlines-\prevgraf \prevgraf\z@
\ifnum\c@WF@wrappedlines<\tw@ \WF@finale \fi}

\makeatother



%---
%makro, co da obrazek doprava a ostatni text ho obteka
%(bez toho predchazejiciho makra to ale poradne nebeha)
%pouziti:
%\obrazekvpravo{adresa}{popisek}{label}{procento sirky}
\long\def\obrazekvpravo#1#2#3#4{

\setlength\intextsep{-20pt}

    \begin{wrapfigure}{r}{#4\textwidth}
      \begin{center}
          \vspace{-10pt}
          
        \includegraphics[width=#4\textwidth]{#1}
        \vspace{-10pt}
        
      \end{center}
      
      \caption{#2}
      \label{#3}
      
      
    \end{wrapfigure}

\setlength\intextsep{0pt}

    
}




%---
%makro pro pripady, kdy wrapfigure neco mrsi
%je to docela pekelne
%je nutne mu dat jak text vpravo, tak text vlevo
%a nevim, jestli bude 100% fungovat, ale doufam, ze jo

%pouziti:
%\obrazekvpravominipage{adresa}{popisek}{label}{procento sirky}{1 - procento sirky}{text vlevo}
\long\def\obrazekvpravominipage#1#2#3#4#5#6{

\noindent\begin{minipage}{#5\linewidth}
\vspace{0pt}
#6
\end{minipage}
\hspace{0.5cm}
\noindent\begin{minipage}{#4\linewidth}
\vspace{0pt}
\centering
\includegraphics[width=0.9\textwidth]{#1}
\captionof{figure}{#2}
\label{#3}
\end{minipage}

}

%KONEC PEKELNYCH MAKER
%=====================

\def\lebka{\includegraphics[width=3mm]{../common/lebka}}

% makra pro poznamku u vyrokove a predikatove logiky
\def\vl{ -- ve v�rokov� logice}
\def\pl{ -- v predik�tov� logice}


%Vacsina prostredi je dvojjazicne. V pripade, ze znenie napr pozorovania je pisane po slovensky, malo by byt po slovensky aj oznacenie.

\def\ifis#1#2{\ifthenelse{\equal{#1}{0}}{}{#2}}

%\newenvironment{e}[3]{\pagebreak[2]\noindent\textbf{\ifis{#2}{$\bigstar$}#1}\ifis{#3}{\emph{~(#3)}}\par\noindent\leftskip 10pt }{\odrovnej\par\bigskip}
\newenvironment{e}[3]{\pagebreak[2]\noindent\textbf{\ifis{#2}{\lebka~}#1}\ifis{#3}{\emph{~(#3)}}\ifis{#2}{~\lebka}\par\noindent\leftskip 10pt }{\odrovnej\par\bigskip}


%obecne prostredie, ktore ma vyuzitie pri specialnych odstavcoch ako (uloha, algoritmus...) aby nevzniklo dalsich x prostredi
\newenvironment{obecne}[1]{\pagebreak[2]\noindent\textbf{#1}\par\noindent\leftskip 10pt}{\odrovnej\par\bigskip}

\newenvironment{report}{\pagebreak[2]\noindent\textbf{Report}\em\par\noindent\leftskip 10pt}{\par\bigskip}

%\newenvironment{reportN}[1]{\pagebreak[2]\noindent\textbf{Report~}\emph{(#1)}\emph\par\noindent\leftskip 10pt}{\odrovnej\par\bigskip}
\newenvironment{reportN}[1]{\pagebreak[2]\noindent\textbf{Report~}\emph{(#1)}\em\par\noindent\leftskip 10pt}{\odrovnej\par\bigskip}

\newenvironment{penumerate}{
\begin{enumerate}
  \setlength{\itemsep}{1pt}
  \setlength{\parskip}{0pt}
  \setlength{\parsep}{0pt}
  %\setlength{\topsep}{200pt}
  \setlength{\partopsep}{200pt}
}{\end{enumerate}}

\def\pismenka{\numberedlistdepth=2} %pouzit, ked clovek chce opismenkovany zoznam...

\newenvironment{pitemize}{
\begin{itemize}
  \setlength{\itemsep}{1pt}
  \setlength{\parskip}{0pt}
  \setlength{\parsep}{0pt}
}{\end{itemize}}

\definecolor{gris}{gray}{0.95}
\newcommand{\ramcek}[2]{\begin{center}\fcolorbox{white}{gris}{\parbox{#1}{#2}}\end{center}\par}


\title{\LARGE U�ebn� texty k st�tn� bakal��sk� zkou�ce \\ Matematika \\ Diskr�tn� matematika}

\begin{document}

\maketitle

\newpage
\setcounter{section}{15}
\def\Nat{\mathbb{N}}
\def\Real{\mathbb{R}}
\def\Pot{\mathcal{P}}

\section{Diskr�tn� matematika}

\begin{e}{Po�adavky}{0}{0}
\begin{pitemize}
\item Uspo��dan� mno�iny
\item Mno�inov� syst�my, p�rov�n�, p�rov�n� v bipartitn�ch grafech (syst�my r�zn�ch reprezentant�)
\item Kombinatorick� po��t�n�
\item Princip inkluze a exkluze
\item Latinsk� �tverce a projektivn� roviny.
\end{pitemize}
\end{e}

\subsection{Uspo��dan� mno�iny}

\begin{e}{Definice}{0}{Kart�zsk� sou�in}
Nech� $X$ a $Y$ jsou mno�iny. Symbolem $X \times Y$ ozna��me mno�inu v�ech uspo��dan�ch dvojic tvaru $(x,y)$, kde $x \in X$ a $y \in Y$. Form�ln� zaps�no:
$X \times Y = \{ (x,~y); x \in X, y \in Y \}$ se naz�v� \emph{kart�zsk� sou�in} mno�in $X$ a $Y$.

Kart�zsk� sou�in $X \times X$ n�kdy zna��me jako $X^2$. 
\end{e}

\begin{e}{Definice}{0}{relace}
Relace R je mno�ina uspo��dan�ch dvojic. Jsou-li $X$ a $Y$ mno�iny, naz�v� se libovoln� podmno�ina kart�zsk�ho sou�inu $X \times Y$ \emph{relac�} mezi $X$ a $Y$. Zdaleka nejd�le�it�j�� p��pad je $X=Y$. V takov�m p��pad� mluv�me o relaci na $X$, co� je tedy libovoln� podmno�ina $X^2$.

N�le��-li $(x,y)$ relaci $R$, tedy $(x,y) \in R$, ��k�me, �e $x$ a $y$ jsou v relaci $R$. Zna��me $xRy$.
\end{e}

\begin{e}{Definice}{0}{druhy relac�}
Relace $X$ m��e b�t:
\begin{pitemize}
        \item \emph{reflexivn�}, jestli�e pro ka�dn� $x \in X$ plat� $xRx$
        \item \emph{ireflexivn�}, jestli�e plat� $xRy\ \Rightarrow\ x\neq y$
        \item \emph{symetrick�}, jestli�e $xRy \Rightarrow yRx$
        \item \emph{tranzitivn�}, jestli�e $xRy \wedge yRz \Rightarrow xRz$
        \item \emph{antisymetrick�}, jestli�e $xRy \wedge yRx \Rightarrow x=y$
\end{pitemize}
\end{e}

\begin{e}{Definice}{0}{ekvivalence}
�ekneme, �e relace $R$ na $X$ je \emph{ekvivalence} na $X$, jestli�e je symetrick�, reflexivn� a tranzitivn�.
\end{e}

\begin{e}{Definice}{0}{uspo��d�n�, uspo��dan� mno�ina}
\emph{Uspo��d�n�} na n�jak� mno�in� $X$ je ka�d� relace na $X$, kter� je reflexivn�, tranzitivn� a antisymetrick�. \emph{Uspo��dan� mno�ina} je dvojice $(X,R)$, kde $X$ je mno�ina a $R$ je uspo��d�n� na $X$.

Pro uspo��d�n� se pou��v� zna�en� $\le$ nebo $\preceq$. Pro ka�d� uspo��d�n� lze odvodit \uv{ostrou nerovnost} $<$ nebo $\prec$, kde plat�, �e $x<y \Leftrightarrow x \le y \wedge x \neq y$.
\end{e}

\begin{e}{P��klady}{0}{0}
Uspo��dan� mno�iny:
\begin{pitemize}
    \item $(\Nat,\leq)$
    \item $(\Nat,|)$, kde \uv{$|$} je relace \uv{d�l�}
    \item $(\Pot(X),\subseteq)$, kde $\Pot(X)$ ozna�uje mno�inu v�ech podmno�in a $\subseteq$ norm�ln� mno�inovou inkluzi
    \item orientovan� acyklick� graf $(V,E)$ s relac� $\rho:(a,b)\in\rho \equiv^{def} \textit{ existuje cesta z }a\textit{ do }b $
\end{pitemize}
\end{e}

\begin{e}{Definice}{0}{line�rn� uspo��d�n�}
\emph{Line�rn� uspo��d�n�} je takov� uspo��d�n�, kde pro ka�d� dva prvky $x$ a $y$ plat� bu� $x \leq y$ nebo $y \leq x$. N�kdy se tak� naz�v� \emph{�pln� uspo��d�n�}.

Uspo��d�n�, kter� nen� �pln�, naz�v�me \emph{��ste�n�m uspo��d�n�m}.
\end{e}

\begin{e}{Definice}{0}{bezprost�edn� p�ech�dce}
\emph{Bezprost�edn� p�edch�dce} prvku $y$ je takov� prvek $x$, pro kter� plat� $x < y$, a neexistuje ��dn� takov� $t$, �e $x < t < y$.
\end{e}

\begin{e}{Pozn�mka}{0}{Hasse�v diagram}
P�i zn�zor�ovan� se uspo��dan� mno�ina zakresluje pomoc� bod� a �ipek, jako kter�koliv jin� relace. Proto�e t�chto �ipek by bylo mnoho, vych�z� se z tranzitivity a zn�zor�uj� se �ipky pouze mezi prvky a jejich bezprost�edn�mi p�edch�dci. P�ijmeme-li nav�c konvenci, �e v obr�zku povedou v�echny �ipky nahoru, nen� t�eba zakreslovat sm�r �ipek, pouze spojnice bod�. Takov�to zn�zorn�n� se pak naz�v� \emph{Hasse�v diagram}.
\end{e}

\begin{e}{Pozn�mka}{0}{uspo��d�n� na kart�zsk�m sou�inu}
M�m-li dv� mno�iny $A$ a $B$ a na nich uspo��d�n� $\leq_A$ a $\leq_B$ m��u definovat \uv{slo�en� uspo��d�n�}
\begin{pitemize}
    \item \uv{po slo�k�ch} -- $(a_1,b_1)\leq(a_2,b_2)\equiv^{\textrm{def}} a_1\leq_A a_2 \wedge b_1 \leq_B b_2$
    \item \uv{lexikograficky} -- $(a_1,b_1)\leq(a_2,b_2)\equiv^{\textrm{def}} a_1\leq_A a_2 \vee ( a_1=a_2 \wedge b_1\leq_B b_2)$
\end{pitemize}
\end{e}

\begin{e}{Definice}{0}{0}
��k�me, �e $(X,R)$ a $(Y,P)$ jsou \emph{isomorfn� uspo��dan� mno�iny}, pokud existuje n�jak� vz�jemn� jednozna�n� zobrazen� $f:X \rightarrow Y$ takov�, �e pro ka�d� $x,y \in X$ plat� $xRy$ pr�v� kdy� $f(x) P f(y)$.
\end{e}

\begin{e}{Definice}{0}{P�eduspo��d�n�}
\emph{P�eduspo��d�n�} nazveme ka�dou relaci, kter� je reflexivn� a transitivn� (nebudeme tedy po�adovat antisymetrii -- mohou vznikat \uv{cykly}).
\end{e}

\begin{e}{Pozn�mka}{0}{0}
M�m-li mno�inu $X$ s relac� $\sim$, kter� je p�eduspo��d�n�, potom relace $\sim$ je uspo��d�n� na mno�in� $X/\rho$ (rozklad podle t��d ekvivalence $\rho$), kde $a\rho b\equiv (a\sim b \wedge b\sim a)$.
\end{e}

\begin{e}{Definice}{0}{0}
\emph{Nez�visl� syst�m $M$ podmno�in mno�iny $X$} je podmno�ina $P(X)$ takov�, �e ka�d� dv� mno�iny $A, B \in M$ jsou neporovnateln� relac� n�le�en�.
\end{e}

\begin{e}{V�ta}{0}{Spernerova}
Libovoln� nez�visl� syst�m podmno�in $n$-prvkov� mno�iny m� nejv�� $\binom{n}{\lfloor\frac{n}{2}\rfloor}$ mno�in.
\end{e}

TODO: Pat�� sem dobr� uspo��d�n� a Zernelova v�ta??? (to by znamenalo p�idat sem i supremum, infimum, �et�zec, nejmen�� a nejv�t�� prvek)

\subsection{Mno�inov� syst�my, p�rov�n�, p�rov�n� v bipartitn�ch grafech (syst�my r�zn�ch reprezentant�)}

\begin{e}{Definice}{0}{0}
Nech� $X$ a $I$ jsou mno�iny. \emph{Mno�inov�m syst�mem nad $X$} nazveme $|I|$-tici $M=\{M_i; i \in I\}$, kde $M_i \subseteq X$

Je tedy mo�n�, aby se tam t� mno�ina objevila v�ckr�t.
\end{e}

\begin{e}{Definice}{0}{syst�m r�zn�ch reprezentant�}
\emph{Syst�m r�zn�ch reprezentant�} (SRR) je funkce $f: I \rightarrow X$ takov�, �e $\forall i \in I: f(i) \in M_i$ a $f$ je prost�.

Jin�mi slovy, SRR je v�b�r jednoho prvku z ka�d� mno�iny $M_i$ tak, �e v�echny vybran� prvky jsou navz�jem r�zn�. Obecn� se tedy neuva�uj� nekone�n� syst�my.
\end{e}

\begin{e}{Definice}{0}{p�rov�n�}
\emph{P�rov�n� v grafu} $G$ je mno�ina hran $F \subseteq E(G)$ takov�, �e ka�d� vrchol grafu $G$ pat�� nejv��e do jedn� hrany z $F$.

Ekvivalentn� definice jsou p�es stupe� vrcholu (ka�d� vrchol m� stupe� nejv��e 1) nebo p�es disjunktnost hran (ka�d� dv� jsou disjunktn� - ��dn� dv� nemaj� spole�n� vrchol).
\end{e}

\begin{e}{Definice}{0}{bipartitn� graf}
\emph{Bipartitn� graf} je takov� graf $G$, kde mno�inu vrchol� $V(G)$ m��eme rozd�lit na dv� disjunktn� podmno�iny $V_1$ a $V_2$ takov�, �e ka�d� hrana z $E(G)$ spojuje v�dy vrchol z $V_1$ s vrcholem z $V_2$.
\end{e}

\begin{e}{Definice}{0}{0}
\emph{Inciden�n�m grafem} mno�inov�ho syst�mu $M$ nad mno�inou $X$ nazveme bipartitn� graf
$$
B_M = \left( I \cup X, \left\{ \{i,x\}, x \in M_i \right\} \right)
$$
\noindent
V podstat� si ka�d� prvek z $X$ i ka�d� index $I$ ozna��me vrcholem a spoj�me ka�d� index $i$ s prvky $x$, kter� n�le�� do $M_i$. Z inciden�n�ho grafu pak lze nahl�dnout existenci SRR - ten existuje, pr�v� kdy� v inciden�n�m grafu existuje p�rov�n� velikosti $|I|$.
\end{e}

\begin{e}{V�ta}{0}{Hallova}
Syst�m r�zn�ch reprezentant� v $M$ existuje, pr�v� kdy� pro ka�dou $J \subseteq I$ je $$\left| \bigcup_{j \in J} M_j\right| \ge |J|$$
\end{e}

\begin{e}{Definice}{0}{perfektn� p�rov�n�}
\emph{Perfektn�m p�rov�n�m} nazveme p�rov�n� $M$ v grafu $G$, pro kter� plat� $$|M| = \frac{|V(G)|}{2}$$

\noindent
Tedy perfektn� p�rov�n� je takov�, kde je ka�d� vrchol sp�rovan� s n�jak�m jin�m vrcholem. Dal��m d�le�it�m pojmem je \emph{maxim�ln� p�rov�n�}, co� je v podstat� nejlep�� mo�n� p�rov�n� (pokr�v� nejv�t�� mo�n� po�et vrchol�), jak�ho jsme v dan�m grafu schopni dos�hnout.
\end{e}


\begin{e}{V�ta}{0}{O p�rov�n� v bipartitn�m grafu}
Bu� $G=(A\cup B,E)$ graf se dv�ma partitami $A$ a $B$ a $E$ nepr�zdn� mno�ina hran. Jestli�e plat� $\deg u \geq \deg v\ \forall u\in A,\forall v\in B$, potom existuje p�rov�n� velikosti $|A|$. D�ky tomu pokud m� bipartitn� graf v�echny vrcholy stejn�ho stupn�, pak m� perfektn� p�rov�n�.

\begin{e}{D�kaz}{0}{0}
P�evede se na Hallovu v�tu s pou�it�m \uv{okol�} vrchol� (tj. bod� p��mo spojen�ch s vrcholem hranou).
\end{e}
\end{e}

\begin{e}{V�ta}{0}{Tutteova}
Graf $(V,E)$ m� perfektn� p�rov�n�, pr�v� kdy� pro ka�dou mno�inu vrchol� $A\subseteq V$ plat�:
$$c_l(G\setminus A)\leq |A|$$
(tj. po�et komponent souvislosti s lich�ch po�tem vrchol� v indukovan�m podgrafu je men�� ne� velikost mno�iny vrchol�). T�to vlastnosti se tak� n�kdy ��k� \emph{Tutteova podm�nka}.
\end{e}

\begin{e}{Algoritmus}{0}{Edmonds�v algoritmus na perfektn� p�rov�n�}
Vstupem algoritmu je graf $G=(V,E)$ a libovoln� p�rov�n� $M$ (i pr�zdn�). V�stupem je p�rov�n� $M'$, kter� alespo� o 1 v�t�� ne� $M$, pokud je n��eho takov�ho mo�n� dos�hnout. Postup v�po�tu:
\begin{penumerate}
    \item Vybudujeme maxim�ln� mo�n� \uv{Edmonds�v les} p�rov�n� $M$ -- do nult� hladiny um�st�me voln� vrcholy a prohled�v�n�m do ���ky sestroj�me max. strom takov�, �e se st��daj� p�rovac� a nep�rovac� hrany (mezi sudou a lichou hladinou jen nep�rovac�, mezi lichou a sudou jen p�rovac�).

N�kter� vrcholy se v lese v�bec neobjev� -- nazveme je \uv{kompost}. Ty jsou u� n�jak sp�rov�ny mezi sebou a nebudeme je pot�ebovat.
    \item Pokud existuje hrana mezi sud�mi hladinami r�zn�ch strom�, m�me \uv{volnou st��davou cestu} (tj. cestu lich� d�lky mezi 2 voln�mi vrcholy, na kter� se st��daj� nep�rovac� a p�rovac� hrany), na n� m��eme zalternovat hrany p�rov�n� a to tak zv�t�it o 1 a skon�it.
    \item Pokud existuje hrana mezi sud�mi hladinami t�ho� stromu, m�me \uv{kv�t} (tj. kru�nici lich� d�lky, na kter� se st��daj� p�rovac� a nep�rovac� hrany). Kv�t m��eme zkontrahovat a algoritmus rekurzivn� pustit na takto z�skan� graf. Pokud dostaneme
    \begin{penumerate}
        \item star� p�rov�n� beze zm�ny, vr�t�me $M'=M$
        \item jin� (v�t��) p�rov�n� $\hat M$, prohod�me p�rov�n� na cest� v $\hat M$ od vrcholu kv�tu v nejvy��� hladin� (kam jsme kv�t zkontrahovali) k voln�mu vrcholu a p�id�me kv�t zp�tky (a p�rov�n� sed� a je v�t�� ne� $M$)
    \end{penumerate}
    \item Nen�-li ��dn� hrana mezi sud�mi hladinami, vydej $M'=M$.
\end{penumerate}
\par\medskip\noindent
Hlavn� algoritmus jen opakuje v��e popsan� krok, dokud vrac� v�t�� p�rov�n� ne� bylo p�edchoz�. Celkov� slo�itost jednoho kroku je $O((m+n)n)$ a pro cel� algoritmus $O((m+n)n^2)$.
\end{e}
\begin{e}{Report}{0}{?}
Pri dvojke stacili definicie a vety. Podstatna bola Hallova veta a este Tutteova veta, ktoru som nevedel. Pan co skusal M bol tiez fer a trpezlivy, rovnako ma nicim prehnanym netrapil. Co som skutocne nevedel som popravde priznal ako pri matematike, tak pri informatike, nastastie toho nebolo privela. 
\end{e}

\begin{e}{Report}{0}{IOI 21.6.2011}
Parovani, perfektni parovani, nezavisly system, system ruznych
reprezentantu
Sepsal jsem definice, SRR a parovani v bip. grafu, SRR a Hallova podminka,
dukaz Hallovy vety jsem mel rozepsany, ale ani ho po mne nechteli, jen na to
zbezne koukli. 
\end{e}
\newpage
\subsection{Kombinatorick� po��t�n�}

\begin{e}{V�ta}{0}{0}
Nech� $N$ je n�jak� $n$-prvkov� mno�ina a $M$ je $m$-prvkov� mno�ina ($m > 0$). Potom po�et v�ech zobrazen� $f: N \rightarrow M$ je $m^n$.
\end{e}

\begin{e}{V�ta}{0}{0}
Libovoln� $n$-prvkov� mno�ina $X$ m� pr�v� $2^n$ podmno�in.
\end{e}

\begin{e}{V�ta}{0}{0}
Nech� $n>0$. Ka�d� $n$-prvkov� mno�ina m� pr�v� $2^{n-1}$ podmno�in sud� velikosti a pr�v� $2^{n-1}$ podmno�in lich� velikosti.
\end{e}

\begin{e}{V�ta}{0}{0}
Pro $m,n \ge 0$ existuje pr�v� $m(m-1)(m-2) \dots (m-n+1)$ prost�ch zobrazen� $n$-prvkov� mno�iny do $m$-prvkov� mno�iny.
\end{e}

\begin{e}{Definice}{0}{0}
Prost� zobrazen� mno�iny $X$ do sebe se naz�v� \emph{permutace mno�iny} $X$. Tato zobrazen� jsou z�rove� na.

Permutace si m��eme p�edstavit jako p�erovn�n� mno�iny - nap�. $\{4 2 1 3\}$ je permutac� mno�iny $\{1 2 3 4\}$. Jin� z�pis permutac� je pomoc� jejich cykl� (\emph{cyklus} v permutaci je po�ad� prvk�, kde za�nu u n�jak�ho prvku, pokra�uji jeho obrazem v permutaci a toto opakuji, dokud nenaraz�m na prvn� prvek). U cykl� zn�zorn�me ka�d� prvek mno�iny jako bod. Z ka�d�ho bodu vych�z� a do ka�d�ho vch�z� pr�v� po jedn� �ipce. �ipka vych�zej�c� z prvn�ho prvku mno�iny ukazuje na prvn� prvek permutace, �ipka z druh�ho prvku mno�iny na druh� prvek permutace atd. Z�pis se pak prov�d� tak, �e ka�dou kru�nici zap�eme po �ad� zvl᚝ (nap�. $p=((1,4,5,2,8)(3)(6,9,7))$ je z�pis permutace $(4 8 3 5 2 9 6 1 7)$.
\end{e}

\begin{e}{V�ta}{0}{Faktori�l}
Po�et permutac� na $n$-prvkov� mno�in� je $n\cdot(n-1)\cdot\dots\cdot 1$. Toto ��slo pojmenujeme \emph{faktori�l} $n$ a zna��me $n!$.
\end{e}

\begin{e}{Definice}{0}{Binomick� koeficienty}
Nech� $n \ge k$ jsou nez�porn� cel� ��sla. \emph{Binomick� koeficient} neboli kombina�n� ��slo $\binom{n}{k}$ (�teme $n$ nad $k$) je funkce prom�nn�ch $n$, $k$, dan� jako $$\binom{n}{k}=\frac{n(n-1)(n-2)\dots(n-k+1)}{1 \cdot 2 \cdot \dots \cdot k}=\frac{\prod^{k-1}_{i=0} (n-i)}{k!}$$ nebo $$\binom{n}{k}=\frac{n!}{k!(n-k)!}$$
\end{e}

\begin{e}{Definice}{0}{0}
Nech� $X$ je mno�ina a $k$ je nez�porn� cel� ��slo. Pak $\binom{X}{k}$ budeme zna�it \emph{ mno�inu v�ech $k$-prvkov�ch podmno�in} mno�iny $X$.
\end{e}

\begin{e}{V�ta}{0}{0}
Pro ka�dou kone�nou mno�inu $X$ je po�et v�ech jej�ch $k$-prvkov�ch podmno�in roven $\binom{|X|}{k}$.
\end{e}

\begin{e}{V�ta}{0}{Vlastnosti kombina�n�ch ��sel}
Plat�:
$$
\binom{n}{k} = \binom{n}{n-k}
$$

$$
\binom{n-1}{k-1} + \binom{n-1}{k} = \binom{n}{k}
$$

\begin{e}{D�kaz}{0}{0}
Prvn� je z�ejm� ze vzorce pro kombina�n� ��sla, druh� se uka�e jednodu�e pro pou�it� komb. ��sel -- m�jme mno�inu $X$ a v n� prvek $a$. Kolik je podmno�in $X$ obsahuj�c�ch, resp. neobsahuj�c�ch $a$?
\end{e}
\end{e}

\begin{e}{D�sledek}{0}{0}
Z druh� vlastnosti plyne vzhled \emph{Pascalova troj�heln�ku}, tedy �e v $n+1$. ��dku jsou v�dy pr�v� binomick� koeficienty $\binom{n}{0},\binom{n}{1},\dots,\binom{n}{n}$.
\end{e}

\begin{e}{V�ta}{0}{O po�tu zp�sob� z�pisu}
Nez�porn� cel� ��slo $m$ lze zapsat jako sou�et $r$ nez�porn�ch s��tanc� pr�v� $\binom{m+r-1}{r-1}$ zp�soby.

\begin{e}{D�kaz}{0}{0}
D�kazem je onen pokus s rozd�lov�n�m $m$ kuli�ek mezi $r$ p�ihr�dek (nebo sp� vkl�d�n� p�ihr�dek mezi kuli�ky v �ad�).
\end{e}
\end{e}

\begin{e}{V�ta}{0}{Binomick� v�ta}
Pro nez�porn� cel� ��slo $n$ a libovoln� $x,y\in\Real$ plat�: $$(x+y)^n=\sum_{k=0}^{n} \binom{n}{k} x^k y^{n-k}$$
\end{e}

\begin{e}{V�ta}{0}{Multinomick� v�ta}
Pro libovoln� ��sla $x_1,\dots,x_m\in\Real$ a $n\in\Nat_0$ plat�: 
$$(x_1+\dots+x_m)^n=\sum_{\genfrac{}{}{0pt}{}{k_1+\dots+k_m=n}{k_1,\dots,k_m\geq 0}}\binom{n}{k_1,\dots,k_n}x_1^{k_1}x_2^{k_2}\cdot\dots\cdot x_m^{k_m}$$
kde ta v�c uprost�ed v tom vzorci je \emph{multinomick� koeficient}, definovan�:
$$\binom{n}{k_1,\dots,k_n}=\frac{n!}{k_1!k_2!\dots k_n!}$$
\end{e}

\subsection{Princip inkluze a exkluze}

\begin{e}{V�ta}{0}{Princip inkluze a exkluze}
Pro ka�d� soubor $A_1, A_2, \dots, A_n$ kone�n�ch mno�in plat�

$$
\left| \bigcup_{i=1}^{m}A_i \right| = \sum_{k=1}^{n} (-1)^{k-1} {\sum_{I \in \binom{\{1,2,\dots,n\}}{k}} {\left| \bigcap_{i \in I} A_i \right|} }
$$
\end{e}

\begin{e}{D�kaz}{0}{0}
Nech� x je libovoln� prvek $A_1 \cup A_2 \cup \dots \cup A_n$. Kolikr�t p�isp�v� x vlevo a kolikr�t vpravo?

\emph{Vlevo}: jednou - trivi�ln�
\emph{Vpravo}: Nech� j ($1 \le j \le n$) ozna�uje po�et mno�in $A_i$, do kter�ch pat�� $x$. M��eme mno�iny p�ejmenovat aby platilo $x~\in~A_1,~A_2,~\dots,~A_j$ a $x~\notin~A_{j+1},~\dots,~A_n$. Prvek se tedy objevuje v pr�niku ka�d� $k$-tice mno�in $A_1,~A_2,~\dots,~A_j$ (a v ��dn�ch jin�ch). Proto�e existuje pr�v� $\binom{j}{k}$ $k$-prvkov�ch podmno�in $j$-prvkov� mno�iny, bude se $x$ objevovat v $\binom{j}{k}$ pr�nic�ch $k$-tic vybran�ch ze v�ech $n$ prvk�. Velikosti $k$-tic jsou p�itom zapo�teny se znam�nkem $(-1)^{k-1}$, tud� $x$ na prav� stran� p�isp�v� veli�inou

\begin{align*}
j - \binom{j}{2} & + \binom{j}{3} - \dots + (-1)^{j-1}\binom{j}{j} =\\
& = 1 - \binom{j}{0} + \binom{j}{1} - \binom{j}{2} \dots + (-1)^{j-1}\binom{j}{j} =\\
& = 1 - \sum_{i=0}^{j} \binom{j}{i} (-1)^{i}  = 1 - (1-1)^{j} = 1
\end{align*}
\begin{flushright}
$\square$
\end{flushright}
\end{e}


\subsection{Latinsk� �tverce a projektivn� roviny}

\begin{e}{Definice}{0}{Projektivn� rovina}
\emph{Kone�n� projektivn� rovina} je mno�inov� syst�m $(B,P)$, kde $B$ je kone�n� mno�ina a $P$ je syst�m podmno�in mno�iny $B$, spl�uj�c�:
\begin{penumerate}
        \item $\forall p \neq p' \in P: |p \cap p'| = 1$
        \item $\forall x \neq y \in B~\exists p \in P: x,y \in p$
        \item $\exists \textit{ 4 body } a,b,c,d \in B: \forall p \in P: |\{a,b,c,d\} \cap p| \leq 2$
\end{penumerate}

Projektivn� rovinu si lze p�edstavit jako mno�inu bod� $B$ a mno�inu \emph{p��mek} $P$ (jak se ostatn� prvky t�chto mno�in naz�vaj�). Pak si lze podm�nky p�edstavit takto:
\begin{penumerate} 
    \item Ka�d� dv� p��mky se prot�naj� pr�v� v jednom bod�. 
    \item Pro ka�d� dva r�zn� body $x$ a $y$ existuje p��mka, kter� jimi proch�z�. 
    \item Existuj� �ty�i body tak, �e ��dn� 3 z nich nele�� na jedn� p��mce (body v obecn� poloze).
\end{penumerate}
\end{e}

\begin{e}{V�ta}{0}{0}
Bu� $(B,P)$ projektivn� rovina. Potom v�echny jej� p��mky maj� stejn� po�et bod�, tedy $\forall p,q \in P: |p| = |q|$
\end{e}

\begin{e}{Definice}{0}{0}
\emph{��d projektivn� roviny} $(B, P)$ je ��slo $|p|-1$, kde $p$ je libovoln� p��mka z $P$ ($p \in P$).
\end{e}

\begin{e}{V�ta}{0}{0}
Nech� $(B,P)$ je projektivn� rovina ��du $n$. Potom plat�, �e ka�d�m bodem proch�z� pr�v� $n+1$ p��mek a $|B|=|P|=n^2+n+1$
\end{e}

\begin{e}{V�ta}{0}{0}
Jestli�e $n=p^2$, kde $p$ je prvo��slo, pak existuje kone�n� projektivn� rovina ��du $n$.
\end{e}

\begin{e}{Definice}{0}{Latinsk� �tverec}
\emph{Latinsk� �tverec} ��du $n$ je matice $A \in \{1, \dots, n\}^{n \times n}, \;\forall i, j \neq j': A_{ij}~\neq~A_{ij'}$ a $A_{ji}~\neq~A_{j'i}.$

V podstat� se jedn� o �tverec $n$ kr�t $n$, kde v ka�d�m ��dku i sloupci jsou vepsan� v�echna ��sla od $1$ do $n$.
\end{e}

\begin{e}{Definice}{0}{0}
M�jme dva latinsk� �tverce $A, B$ stejn�ch rozm�r� $n \times n$. Pak �ekneme, �e je \emph{ortogon�ln�} (zna��me $A \bot B$), jestli�e plat�:
$$\forall a,b \in \{1,\dots,n\} \;\exists i,j: A_{ij}=a, B_{ij}=b$$

Pokud tedy ty dva latinsk� �tverce \uv{polo��me p�es sebe}, vznikne n�m na ka�d� pozici dvojice ��sel od $1$ do $n$ s t�m, �e ka�d� dvojice je unik�tn�.
\end{e}

\begin{e}{V�ta}{0}{0}
Nech� $M$ je mno�ina latinsk�ch �tverc� ��du $n$, z nich� ka�d� dva jsou navz�jem ortogon�ln�. Potom $|M| \le n-1$
\end{e}

\begin{e}{V�ta}{0}{0}
Pro $n \ge 2$, projektivn� rovina ��du $n$ existuje pr�v� tehdy, kdy� existuje soubor $n-1$ vz�jemn� ortogon�ln�ch latinsk�ch �tverc� ��du $n$.
\par\medskip
\begin{e}{D�kaz}{0}{0}
\begin{enumerate}
    \item \emph{Implikace $\Leftarrow$:} Vezmu jednu -- \uv{vn�j��} -- p��mku projektivn� roviny. Na n� le�� $n+1$ bod�, kter� nazvu $A_0,\dots,A_n$. P��mky jdouc� z krajn�ch bod� $A_0$ a $A_n$ tvo�� m��ku (nazvu ji $T$) -- prot�naj� se v $n^2$ bodech. 

Potom ka�d� z vnit�n�ch bod� $A_i\ (1 ? i ? n-1)$ definuje lat. �tverec: ka�d� p��mka jdouc� z n�jak�ho z vnit�n�ch bod� $A_i$ se s p��mkami z $A_0$ a z $A_n$ protne pr�v� jednou a ka�d� 2 p��mky z $A_i$ proch�z� mimo vn�j�� p��mku r�zn�mi body. Ka�d� p��mka proch. ka�d�m ��dkem i sloupcem m��ky $T$ pr�v� jednou $\Rightarrow$ dost�v�m latinsk� �tverec: 
$$(L^k)_{ij} = l \Leftrightarrow T_{ij} \in p^k_l$$
kde $L^k$ zna�� $k$-t� lat. �tverec, $T_{ij}$ bod m��ky $T$ na sou�adnic�ch $(i,j)$ a $p^k_l$ $l$-tou p��mku jdouc� z bodu $A_k$. 

�tverce jsou ortogon�ln� - sporem nech� pro maj� dva �tverce ($k$-t� a $k'$-t�) na sou�adnic�ch stejnou uspo��danou dvojici hodnot $(a,b)$ na dvou r�zn�ch m�stech. Pak by se p��mky $p^k_a$ a $p^{k'}_b$ prot�naly ve 2 bodech.

    \item \emph{Implikace $\Rightarrow$:} Vytvo��m p��mku $q$ s body $A_0,\dots,A_n$ a m��ku $T$ o $n\times n$ bodech. Do n� p�id�m p��mky $p_{1,1},\dots,p_{n-1,n}$: 
$$p_l^k = \{A_k\} \cup \{T_{ij} | L^k_{ij} = l \}$$ 
Pak je t�eba ov��it axiomy projektivn� roviny.�
\end{enumerate}
\end{e}
\end{e}


\end{document}
